\documentclass[11pt,a4paper]{article}
\usepackage{fontspec,amsmath,amssymb,amsthm,microtype}
\setmainfont{lmroman10-regular.otf}[BoldFont=lmroman10-bold.otf,ItalicFont=lmroman10-italic.otf,BoldItalicFont=lmroman10-bolditalic.otf]
\usepackage{xcolor}
\usepackage[most]{tcolorbox}
\usepackage{mathrsfs}
\providecommand{\lt}{<}
\providecommand{\gt}{>}
\setmonofont{DejaVuSansMono.ttf}[BoldFont=DejaVuSansMono-Bold.ttf,ItalicFont=DejaVuSansMono-Oblique.ttf,BoldItalicFont=DejaVuSansMono-BoldOblique.ttf,Scale=0.86]
\newfontfamily\leanfallback{STIXTwoMath-Regular.otf}[Scale=0.86]
\definecolor{leanrule}{gray}{0.62}
\definecolor{leanback}{gray}{0.962}
\definecolor{leantitle}{gray}{0.22}
\definecolor{leanmuted}{gray}{0.38}
\newtcolorbox{leanbox}[1]{enhanced,breakable,colback=leanback,colframe=leanrule,
  boxrule=0.4pt,arc=0.6mm,left=2.4mm,right=2.4mm,top=1.8mm,bottom=1.8mm,
  before skip=0.8em,after skip=1em,fonttitle=\small\bfseries,coltitle=leantitle,
  colbacktitle=leanback,titlerule=0.3pt,title={#1}}
\newlength{\leancharwidth}
\newenvironment{leancode}{\par\smallskip\begingroup\ttfamily\footnotesize\raggedright
  \settowidth{\leancharwidth}{M}\setlength{\parskip}{0pt}\setlength{\parindent}{0pt}
  \frenchspacing}{\par\endgroup\smallskip}
\newcommand{\leanline}[2]{\hangindent=\dimexpr#1\leancharwidth+4\leancharwidth\relax
  \hangafter=1\noindent\hspace*{#1\leancharwidth}#2\par}
\newcommand{\leansource}[1]{\par{\footnotesize\color{leanmuted}#1}\par}
\newcommand{\leanlabel}[1]{\par\smallskip\noindent{\small\bfseries #1}\par}
\usepackage[margin=25mm]{geometry}
\usepackage[unicode,hidelinks]{hyperref}
\hypersetup{pdfauthor={LangevinLean Project}}
\setlength{\emergencystretch}{3em}
\setlength{\parskip}{0.4em}
\setlength{\parindent}{0pt}
\allowdisplaybreaks[2]
\raggedbottom
\title{\texorpdfstring{A genuine integral estimate for the local defect}{A genuine integral estimate for the local defect}}
\author{LangevinLean Project}
\date{}
\begin{document}\maketitle
Draft mathematical exposition; these notes may contain mistakes.


Let \(H\) be a complete real normed space and \(J:[0,h]\to H\) a square-integrable increment. The local defect is the actual Bochner integral \(D=\int_0^hJ(s)\,ds\). Time Cauchy–Schwarz yields the source constant \(3\) from the primitive square-increment constant \(6\).


\section{\texorpdfstring{Scalar Cauchy–Schwarz through the square-integrable Hilbert space}{Scalar Cauchy–Schwarz through the square-integrable Hilbert space}}


\textbf{Claim.} If \(f,g\) are square-integrable real functions on a measure space, then

\[\int fg\le\left(\int f^2\right)^{1/2}\left(\int g^2\right)^{1/2}.\]


\textbf{Proof.} The squared norm of the equivalence class of \(f\) in \(L^2\) is \(\int f^2\), and likewise for \(g\). Their real inner product is \(\int fg\). Apply the Hilbert-space bound for the inner product and replace each norm by the square root of its squared integral. The identifications hold almost everywhere and therefore preserve the integrals.

\begin{leanbox}{Lean: \texttt{toLp\_\allowbreak{}norm\_\allowbreak{}sq}, \texttt{real\_\allowbreak{}cauchy\_\allowbreak{}schwarz}}
\leansource{\texttt{toLp\_\allowbreak{}norm\_\allowbreak{}sq} \quad WassersteinMixing/\allowbreak{}LocalError.lean, lines 19–25}
\begin{leancode}
\leanline{0}{private\ theorem\ toLp\_norm\_sq\ \{Ω\ :\ Type*\}\ [MeasurableSpace\ Ω]\ (Q\ :\ Measure\ Ω)}
\leanline{4}{(f\ :\ Ω\ →\ ℝ)\ (hf\ :\ MemLp\ f\ 2\ Q)\ :}
\leanline{4}{‖hf.toLp\ f‖\ \textasciicircum{}\ 2\ =\ ∫\ ω,\ f\ ω\ \textasciicircum{}\ 2\ ∂Q}
\end{leancode}
\leansource{\texttt{real\_\allowbreak{}cauchy\_\allowbreak{}schwarz} \quad WassersteinMixing/\allowbreak{}LocalError.lean, lines 27–40}
\begin{leancode}
\leanline{0}{private\ theorem\ real\_cauchy\_schwarz\ \{Ω\ :\ Type*\}\ [MeasurableSpace\ Ω]\ (Q\ :\ Measure\ Ω)}
\leanline{4}{(f\ g\ :\ Ω\ →\ ℝ)\ (hf\ :\ MemLp\ f\ 2\ Q)\ (hg\ :\ MemLp\ g\ 2\ Q)\ :}
\leanline{4}{(∫\ ω,\ f\ ω\ *\ g\ ω\ ∂Q)\ ≤}
\leanline{6}{Real.sqrt\ (∫\ ω,\ f\ ω\ \textasciicircum{}\ 2\ ∂Q)\ *\ Real.sqrt\ (∫\ ω,\ g\ ω\ \textasciicircum{}\ 2\ ∂Q)}
\end{leancode}
\leanlabel{How the Lean proof follows the argument}
\texttt{to\allowbreak{}Lp\_\allowbreak{}norm\_\allowbreak{}sq} uses the real \texttt{L2.\allowbreak{}inner\_\allowbreak{}def} and the almost-everywhere representative equality. \texttt{real\_\allowbreak{}cauchy\_\allowbreak{}schwarz} identifies the integral product with the genuine inner product and applies \texttt{real\_\allowbreak{}inner\_\allowbreak{}le\_\allowbreak{}norm}.
\end{leanbox}

\section{\texorpdfstring{Time Cauchy–Schwarz for the Bochner integral}{Time Cauchy–Schwarz for the Bochner integral}}


\textbf{Claim.} For \(h\ge0\) and square-integrable \(J\) on \((0,h]\),

\[\left\|\int_0^hJ(s)\,ds\right\|^2\le h\int_0^h\|J(s)\|^2\,ds.\]


\textbf{Proof.} The norm of the Bochner integral is at most \(\int_0^h\|J(s)\|\,ds\). Apply scalar Cauchy–Schwarz to \(\|J\|\) and the constant function \(1\). The latter has squared integral \(h\). Both resulting sides are nonnegative, so squaring gives the stated inequality. This also defines the local defect \(D\) as the genuine integral, with all square-integrability hypotheses retained.

\begin{leanbox}{Lean: \texttt{time\_\allowbreak{}integral\_\allowbreak{}norm\_\allowbreak{}sq}, \texttt{stationaryLocalDefect}}
\leansource{\texttt{time\_\allowbreak{}integral\_\allowbreak{}norm\_\allowbreak{}sq} \quad WassersteinMixing/\allowbreak{}LocalError.lean, lines 45–63}
\begin{leancode}
\leanline{0}{theorem\ time\_integral\_norm\_sq\ (increment\ :\ ℝ\ →\ H)\ (h\ :\ ℝ)\ (hh\ :\ 0\ ≤\ h)}
\leanline{4}{(hi\ :\ MemLp\ increment\ 2\ (volume.restrict\ (Ioc\ 0\ h)))\ :}
\leanline{4}{‖∫\ s\ in\ (0\ :\ ℝ)..h,\ increment\ s‖\ \textasciicircum{}\ 2\ ≤}
\leanline{6}{h\ *\ ∫\ s\ in\ (0\ :\ ℝ)..h,\ ‖increment\ s‖\ \textasciicircum{}\ 2}
\end{leancode}
\leansource{\texttt{stationaryLocalDefect} \quad WassersteinMixing/\allowbreak{}LocalError.lean, lines 66–67}
\begin{leancode}
\leanline{0}{def\ stationaryLocalDefect\ (increment\ :\ ℝ\ →\ H)\ (h\ :\ ℝ)\ :\ H\ :=}
\leanline{2}{∫\ s\ in\ (0\ :\ ℝ)..h,\ increment\ s}
\end{leancode}
\leanlabel{How the Lean proof follows the argument}
\texttt{time\_\allowbreak{}integral\_\allowbreak{}norm\_\allowbreak{}sq} combines \texttt{interval\allowbreak{}Integral.\allowbreak{}norm\_\allowbreak{}integral\_\allowbreak{}le\_\allowbreak{}integral\_\allowbreak{}norm} with the preceding scalar bound for the restricted interval measure. It evaluates the measure of the interval and uses the nonnegative square-root identities before squaring. \texttt{stationary\allowbreak{}Local\allowbreak{}Defect} is the Bochner interval integral itself.
\end{leanbox}

\section{\texorpdfstring{The exact local defect constant and the displacement polynomial}{The exact local defect constant and the displacement polynomial}}


\textbf{Claim.} If \(\|J(s)\|^2\le6s\beta^2d\) on \((0,h]\), then

\[\|D\|^2\le3h^3\beta^2d.\]


\textbf{Proof.} Square integrability makes \(\|J\|^2\) integrable. Integrate its upper bound and use the preceding time estimate:

\[\|D\|^2\le h\int_0^h6s\beta^2d\,ds
=h\cdot3h^2\beta^2d.\]


For the auxiliary source displacement polynomial, \(0\le s\le h\), \(\beta\ge0\) and \(h\beta\le1\) imply \(s\beta\le1\), hence \(2s^2\beta+4s\le6s\). This last simplification supplies no independent probabilistic displacement estimate.

\begin{leanbox}{Lean: \texttt{stationaryLocalDefect\_\allowbreak{}energy\_\allowbreak{}le}, \texttt{displacement\_\allowbreak{}polynomial\_\allowbreak{}le}}
\leansource{\texttt{stationaryLocalDefect\_\allowbreak{}energy\_\allowbreak{}le} \quad WassersteinMixing/\allowbreak{}LocalError.lean, lines 70–89}
\begin{leancode}
\leanline{0}{theorem\ stationaryLocalDefect\_energy\_le\ (increment\ :\ ℝ\ →\ H)\ (h\ β\ d\ :\ ℝ)}
\leanline{4}{(hh\ :\ 0\ ≤\ h)\ (hi\ :\ MemLp\ increment\ 2\ (volume.restrict\ (Ioc\ 0\ h)))}
\leanline{4}{(hbound\ :\ ∀\ s\ ∈\ Ioc\ 0\ h,\ ‖increment\ s‖\ \textasciicircum{}\ 2\ ≤\ 6\ *\ s\ *\ β\ \textasciicircum{}\ 2\ *\ d)\ :}
\leanline{4}{‖stationaryLocalDefect\ increment\ h‖\ \textasciicircum{}\ 2\ ≤\ 3\ *\ h\ \textasciicircum{}\ 3\ *\ β\ \textasciicircum{}\ 2\ *\ d}
\end{leancode}
\leansource{\texttt{displacement\_\allowbreak{}polynomial\_\allowbreak{}le} \quad WassersteinMixing/\allowbreak{}LocalError.lean, lines 92–96}
\begin{leancode}
\leanline{0}{theorem\ displacement\_polynomial\_le\ (s\ h\ β\ :\ ℝ)\ (hs\ :\ 0\ ≤\ s)\ (hsh\ :\ s\ ≤\ h)}
\leanline{4}{(hβ\ :\ 0\ ≤\ β)\ (hstep\ :\ h\ *\ β\ ≤\ 1)\ :}
\leanline{4}{2\ *\ s\ \textasciicircum{}\ 2\ *\ β\ +\ 4\ *\ s\ ≤\ 6\ *\ s}
\end{leancode}
\leanlabel{How the Lean proof follows the argument}
\texttt{stationary\allowbreak{}Local\allowbreak{}Defect\_\allowbreak{}energy\_\allowbreak{}le} proves integrability of the squared norm from \texttt{Mem\allowbreak{}Lp}, applies monotonicity of the interval integral and evaluates the linear polynomial exactly. \texttt{displacement\_\allowbreak{}polynomial\_\allowbreak{}le} isolates the arithmetic consequence of the stable-step restriction.
\end{leanbox}

\end{document}
